\section{AIME II 2005}
        \begin{enumerate}
        	\item (\textit{AIME II 2005}) A game uses a deck of $n$ different cards, where $n$ is an integer and $n \geq 6.$ The number of possible sets of 6 cards that can be drawn from the deck is 6 times the number of possible sets of 3 cards that can be drawn. Find $n.$


 
 \textbf{Answer: }013



\textbf{Solution}

The number of ways to draw six cards from $n$ is given by the \href{https://artofproblemsolving.com/wiki/index.php?title=Binomial\_coefficient}{binomial coefficient} ${n \choose 6} = \frac{n\cdot(n-1)\cdot(n-2)\cdot(n-3)\cdot(n-4)\cdot(n-5)}{6\cdot5\cdot4\cdot3\cdot2\cdot1}$.  

 The number of ways to choose three cards from $n$ is ${n\choose 3} = \frac{n\cdot(n-1)\cdot(n-2)}{3\cdot2\cdot1}$.

 We are given that ${n\choose 6} = 6 {n \choose 3}$, so $\frac{n\cdot(n-1)\cdot(n-2)\cdot(n-3)\cdot(n-4)\cdot(n-5)}{6\cdot5\cdot4\cdot3\cdot2\cdot1} = 6 \frac{n\cdot(n-1)\cdot(n-2)}{3\cdot2\cdot1}$.

 Cancelling like terms, we get $(n - 3)(n - 4)(n - 5) = 720$.

 We must find a \href{https://artofproblemsolving.com/wiki/index.php?title=Factoring}{factorization} of the left-hand side of this equation into three consecutive \href{https://artofproblemsolving.com/wiki/index.php?title=Integer}{integers}. Since 720 is close to $9^3=729$, we try 8, 9, and 10, which works, so $n - 3 = 10$ and $n = \boxed{013}$.

 




	\item (\textit{AIME II 2005}) A hotel packed breakfast for each of three guests. Each breakfast should have consisted of three types of rolls, one each of nut, cheese, and fruit rolls. The preparer wrapped each of the nine rolls and once wrapped, the rolls were indistinguishable from one another. She then randomly put three rolls in a bag for each of the guests. Given that the probability each guest got one roll of each type is $\frac mn,$ where $m$ and $n$ are relatively prime integers, find $m+n.$


 
 \textbf{Answer: }079



\textbf{Solution 1}

Use \href{https://artofproblemsolving.com/wiki/index.php?title=Construction}{construction}. We need only calculate the probability the first and second person all get a roll of each type, since then the rolls for the third person are determined. 

 
                \begin{itemize}
                    \item Person 1: $\frac{9 \cdot 6 \cdot 3}{9 \cdot 8 \cdot 7} = \frac{9}{28}$
                \end{itemize}


                
                \begin{itemize}
                    \item Person 2: $\frac{6 \cdot 4 \cdot 2}{6 \cdot 5 \cdot 4} = \frac 25$\item Person 3: One roll of each type is left, so the probability here is $1$.
                \end{itemize}


                Our answer is thus $\frac{9}{28} \cdot \frac{2}{5} = \frac{9}{70}$, and $m + n = \boxed{79}$.

 





\textbf{Solution 2}

Call the three different types of rolls as A, B, and C. We need to arrange 3As, 3Bs, and 3Cs in a string such that A, B, and C appear in the first three, second three, and the third three like ABCABCABC or BCABACCAB. This can occur in $\left(\frac{3!}{1!1!1!}\right)^3 = 6^3 = 216$ different manners. The total number of possible strings is $\frac{9!}{3!3!3!} = 1680$. The solution is therefore $\frac{216}{1680} = \frac{9}{70}$, and $m + n = \boxed{79}$.

 





\textbf{Solution 3}

The denominator of m/n is equal to the total amount of possible roll configurations given to the three people. This is equal to ${9 \choose 3}{6 \choose3}$ as the amount of ways to select three rolls out of 9 to give to the first person is ${9 \choose 3}$, and three rolls out of 6 is ${6 \choose3}$. After that, the three remaining rolls have no more configurations. 

 The numerator is the amount of ways to give one roll of each type to each of the three people, which can be done by defining the three types of rolls as x flavored, y flavored, and z flavored. 

 xxx, yyy, zzz

 So you have to choose one x, one y, and one z to give to the first person. There are 3 xs, 3 ys, and 3 zs to select from, giving $3^3$ combinations. Multiply that by the combinations of xs, ys, and zs for the second person, which is evidently $2^3$ since there are two of each letter left. 

 $(27*8)/{9 \choose 3}{6 \choose3}$ simplifies down to our fraction m/n, which is $9/70$. Adding them up gives $9 + 70 = \boxed{79}$.

 





\textbf{Solution 4}

Let the objects be $A_1$, $A_2$, $A_3$ with similar labeling for $B$ and $C$. There are $9!$ ways to order these nine objects in a line.

 We then label three "zones" corresponding to the three people's meals as the first, second, and third objects in the line corresponding to the first zone, the fourth, fifth, and sixth to the second zone, and the remainder to the third zone. Each zone must contain one of each $A$, $B$, and $C$. There are $6$ ways to place the three $A$ rolls in three different zones, and the same goes for $B$ and $C$. In each zone, there are $6$ ways to order the three rolls, so there are in total $6^6$ legal orderings.

 Thus the desired probability is $\frac{6^6}{9!}=\frac{9}{70}$, so the answer is $9+70=\boxed{79}$.

 




	\item (\textit{AIME II 2005}) An infinite geometric series has sum 2005. A new series, obtained by squaring each term of the original series, has 10 times the sum of the original series. The common ratio of the original series is $\frac mn$ where $m$ and $n$ are relatively prime integers. Find $m+n.$


 
 \textbf{Answer: }802



\textbf{Solution 1}

Let's call the first term of the original \href{https://artofproblemsolving.com/wiki/index.php?title=Geometric\_series}{geometric series} $a$ and the common ratio $r$, so $2005 = a + ar + ar^2 + \ldots$.  Using the sum formula for \href{https://artofproblemsolving.com/wiki/index.php?title=Infinite}{infinite} geometric series, we have $\;\;\frac a{1 -r} = 2005$.  Then we form a new series, $a^2 + a^2 r^2 + a^2 r^4 + \ldots$.  We know this series has sum  $20050 = \frac{a^2}{1 - r^2}$.  Dividing this equation by $\frac{a}{1-r}$, we get $10 = \frac a{1 + r}$.  Then $a = 2005 - 2005r$ and $a = 10 + 10r$ so $2005 - 2005r = 10 + 10r$, $1995 = 2015r$ and finally $r = \frac{1995}{2015} = \frac{399}{403}$, so the answer is $399 + 403 = \boxed{802}$.  

 (We know this last fraction is fully reduced by the \href{https://artofproblemsolving.com/wiki/index.php?title=Euclidean\_algorithm}{Euclidean algorithm} -- because $4 = 403 - 399$, $\gcd(403, 399) | 4$.  But 403 is \href{https://artofproblemsolving.com/wiki/index.php?title=Odd\_integer}{ odd}, so $\gcd(403, 399) = 1$.)

 





\textbf{Solution 2}

We can write the sum of the original series as $a + a\left(\dfrac{m}{n}\right) + a\left(\dfrac{m}{n}\right)^2 + \ldots = 2005$, where the common ratio is equal to $\dfrac{m}{n}$. We can also write the sum of the second series as $a^2 + a^2\left(\dfrac{m}{n}\right)^2 + a^2\left(\left(\dfrac{m}{n}\right)^2\right)^2 + \ldots = 20050$. Using the formula for the sum of an infinite geometric series $S=\dfrac{a}{1-r}$, where $S$ is the sum of the sequence, $a$ is the first term of the sequence, and $r$ is the ratio of the sequence, the sum of the original series can be written as $\dfrac{a}{1-\frac{m}{n}}=\dfrac{a}{\frac{n-m}{n}}=\dfrac{a \cdot n}{n-m}=2005\;\text{(1)}$, and the second sequence can be written as $\dfrac{a^2}{1-\frac{m^2}{n^2}}=\dfrac{a^2}{\frac{n^2-m^2}{n^2}}=\dfrac{a^2\cdot n^2}{(n+m)(n-m)}=20050\;\text{(2)}$. Dividing $\text{(2)}$ by $\text{(1)}$, we obtain $\dfrac{a\cdot n}{m+n}=10$, which can also be written as $a\cdot n=10(m+n)$. Substitute this value for $a\cdot n$ back into $\text{(1)}$, we obtain $10\cdot \dfrac{n+m}{n-m}=2005$. Dividing both sides by 10 yields $\dfrac{n+m}{n-m}=\dfrac{401}{2}$ we can now write a system of equations with $n+m=401$ and $n-m=2$, but this does not output integer solutions. However, we can also write $\dfrac{n+m}{n-m}=\dfrac{401}{2}$ as $\dfrac{n+m}{n-m}=\dfrac{802}{4}$. This gives the system of equations $m+n=802$ and $n-m=4$, which does have integer solutions. Our answer is therefore $m+n=\boxed{802}$ (Solving for $m$ and $n$ gives us $399$ and $403$, respectively, which are co-prime).

 




	\item (\textit{AIME II 2005}) Find the number of positive integers that are divisors of at least one of $10^{10},15^7,18^{11}.$


 
 \textbf{Answer: }435



\textbf{Solution 1}

$10^{10} = 2^{10}\cdot 5^{10}$ so $10^{10}$ has $11\cdot11 = 121$ \href{https://artofproblemsolving.com/wiki/index.php?title=Divisor}{divisors}.

 $15^7 = 3^7\cdot5^7$ so $15^7$ has $8\cdot8 = 64$ divisors.

 $18^{11} = 2^{11}\cdot3^{22}$ so $18^{11}$ has $12\cdot23 = 276$ divisors.

 Now, we use the \href{https://artofproblemsolving.com/wiki/index.php?title=Principle\_of\_Inclusion-Exclusion}{Principle of Inclusion-Exclusion}.  We have $121 + 64 + 276$ total potential divisors so far, but we've overcounted those factors which divide two or more of our three numbers.  Thus, we must subtract off the divisors of their pair-wise \href{https://artofproblemsolving.com/wiki/index.php?title=Greatest\_common\_divisor}{greatest common divisors}.

 $\gcd(10^{10},15^7) = 5^7$ which has 8 divisors.

 $\gcd(15^7, 18^{11}) = 3^7$ which has 8 divisors.

 $\gcd(18^{11}, 10^{10}) = 2^{10}$ which has 11 divisors.

 So now we have $121 + 64 + 276 - 8 -8 -11$ potential divisors.  However, we've now undercounted those factors which divide all three of our numbers.  Luckily, we see that the only such factor is 1, so we must add 1 to our previous sum to get an answer of $121 + 64 + 276 - 8 - 8 - 11 + 1 = \boxed{435}$.

 





\textbf{Solution 2}

We can rewrite the three numbers as $10^{10} = 2^{10}\cdot 5^{10}$, $15^7 = 3^7\cdot5^7$, and $18^{11} = 2^{11}\cdot3^{22}$.Assume that $n$ (a positive integer) is a divisor of one of the numbers. Therefore, $n$ can be expressed as ${p_1}^{e_1}$ or as ${p_2}^{e_2}{p_3}^{e_3}$ where $p_1$, $p_2$ are in $\{2,3,5\}$ and $e_1$, $e_2$ are positive integers.

 If $n$ is the power of a single prime, then there are 11 possibilities ($2^1$ to $2^{11}$) for $p_1=2$, 22 possibilities ($3^1$ to $3^{22}$) for $p_1=3$, 10 possibilities ($5^1$ to $5^{10}$) for $p_1=5$, and 1 possibility if $n=1$. From this case, there are $11+22+10+1=44$ possibilities.

 If $n$ is the product of the powers of two primes, then we can just multiply the exponents of each rewritten product to get the number of possibilities, since each exponent of the product must be greater than 0. From this case, there are $10*10+11*22+7*7=100+242+49=391$ possibilities.

 Adding up the two cases, there are $44+391=\boxed{435}$ positive integers.

 -alpha\_2

 




	\item (\textit{AIME II 2005}) Determine the number of ordered pairs $(a,b)$ of integers such that $\log_a b + 6\log_b a=5, 2 \leq a \leq 2005,$ and $2 \leq b \leq 2005.$


 
 \textbf{Answer: }054



\textbf{Solution 1}

The equation can be rewritten as $\frac{\log b}{\log a} + 6 \frac{\log a}{\log b} = \frac{(\log b)^2+6(\log a)^2}{\log a  \log b}=5$ Multiplying through by $\log a \log b$ and factoring yields $(\log b - 3\log a)(\log b - 2\log a)=0$. Therefore, $\log b=3\log a$ or $\log b=2\log a$, so either $b=a^3$ or $b=a^2$. 

 
                \begin{itemize}
                    \item For the case $b=a^2$, note that $44^2=1936$ and $45^2=2025$. Thus, all values of $a$ from $2$ to $44$ will work.\item For the case $b=a^3$, note that $12^3=1728$ while $13^3=2197$. Therefore, for this case, all values of $a$ from $2$ to $12$ work.
                \end{itemize}


                There are $44-2+1=43$ possibilities for the square case and $12-2+1=11$ possibilities for the cube case. Thus, the answer is $43+11= \boxed{054}$.

 Note that Inclusion-Exclusion does not need to be used, as the problem is asking for ordered pairs $(a,b)$, and not for the number of possible values of $b$. Were the problem to ask for the number of possible values of $b$, the values of $b^6$ under $2005$ would have to be subtracted, which would just be $2$ values:  $2^6$ and $3^6$. However, the ordered pairs where b is to the sixth power are distinct, so they are not redundant. (For example, the pairs (4, 64) and (8, 64).)

 





\textbf{Solution 2}

Let $k=\log_a b$. Then our equation becomes $k+\frac{6}{k}=5$. Multiplying through by $k$ and solving the quadratic gives us $k=2$ or $k=3$. Hence $a^2=b$ or $a^3=b$.  

 For the first case $a^2=b$, $a$ can range from 2 to 44, a total of 43 values.For the second case $a^3=b$, $a$ can range from 2 to 12, a total of 11 values.

 
Thus the total number of possible values is $43+11=\boxed{54}$.

 


 





\textbf{Solution 3(similar to solution 2)}

Using the change of base formula on the second equation to change to base $a$, we get $\log_a(b) + \frac{6 \log_a(a)}{\log_a(b)}$. If we substitute $x$ for $\log_a(b)$, we get $x + \frac{6}{x}$. Multiplying by $x$ on both sides and solving, we get $x=3,2$. Substituting back in, we get $\log_a(b) = 3,2$. That means $a^3 = b$ or $a^2 = b$. Since $b \leq 2005$, we can see that for the cubed case, the maximum $a$ can be without exceeding 2005 is 12(because $13^3 = 2197$) and for the squared case it can be a maximum of 44. Since $a \neq 1$, the number of values is $(44-1)+(12-1) = \boxed{54}$. 

 




	\item (\textit{AIME II 2005}) The cards in a stack of $2n$ cards are numbered consecutively from 1 through $2n$ from top to bottom. The top $n$ cards are removed, kept in order, and form pile $A.$ The remaining cards form pile $B.$ The cards are then restacked by taking cards alternately from the tops of pile $B$ and $A,$ respectively. In this process, card number $(n+1)$ becomes the bottom card of the new stack, card number 1 is on top of this card, and so on, until piles $A$ and $B$ are exhausted. If, after the restacking process, at least one card from each pile occupies the same position that it occupied in the original stack, the stack is named magical. Find the number of cards in the magical stack in which card number 131 retains its original position.


 
 \textbf{Answer: }392



\textbf{Solution 1}

Since a card from B is placed on the bottom of the new stack, notice that cards from pile B will be marked as an even number in the new pile, while cards from pile A will be marked as odd in the new pile. Since 131 is odd and retains its original position in the stack, it must be in pile A. Also to retain its original position, exactly $131 - 1 = 130$ numbers must be in front of it. There are $\frac{130}{2} = 65$ cards from each of piles A, B in front of card 131. This suggests that $n = 131 + 65 = 196$; the total number of cards is $196 \cdot 2 = \boxed{392}$.

 





\textbf{Solution 2}

If you index the final stack $1,2,\dots,2n$, you notice that pile A resides only in the odd indices and has maintained its original order aside from flipping over. The same has happened to pile B except replace odd with even. Thus, if 131 is still at index 131, an odd number, then 131 must be from pile A. The numbers in pile A are the consecutive integers $1,2,\dots, n$. This all leads us to the following equation.
 \[131=2n-2(131)+1\implies2n=\boxed{392}\]

 




	\item (\textit{AIME II 2005}) Let $x=\frac{4}{(\sqrt{5}+1)(\sqrt[4]{5}+1)(\sqrt[8]{5}+1)(\sqrt[16]{5}+1)}.$ Find $(x+1)^{48}.$


 
 \textbf{Answer: }125



\textbf{Solution 1}

We note that in general,

 <center>${} (\sqrt[2^n]{5} + 1)(\sqrt[2^n]{5} - 1) = (\sqrt[2^n]{5})^2 - 1^2 = \sqrt[2^{n-1}]{5} - 1$.

 </center>It now becomes apparent that if we multiply the \href{https://artofproblemsolving.com/wiki/index.php?title=Numerator}{numerator} and \href{https://artofproblemsolving.com/wiki/index.php?title=Denominator}{denominator} of $\frac{4}{ (\sqrt{5}+1) (\sqrt[4]{5}+1) (\sqrt[8]{5}+1) (\sqrt[16]{5}+1) }$ by $(\sqrt[16]{5} - 1)$, the denominator will \href{https://artofproblemsolving.com/wiki/index.php?title=Telescope}{telescope} to $\sqrt[1]{5} - 1 = 4$, so 

 <center>$x = \frac{4(\sqrt[16]{5} - 1)}{4} = \sqrt[16]{5} - 1$.

 </center>It follows that $(x + 1)^{48} = (\sqrt[16]5)^{48} = 5^3 = \boxed{125}$.

 





\textbf{Solution 2 (Clever)}

First, multiply the denominator $\sqrt{5} + 1$ term by its conjugate, and obtain 

 
 \[\frac{\sqrt{5}-1}{(\sqrt[4]{5}+1)(\sqrt[8]{5}+1)(\sqrt[16]{5}+1)}\]

 Define $a = \sqrt[16]{5}$. Then, we have

 
 \[\frac{a^8 - 1}{(a^4 + 1)(a^2+1)(a+1)}\]
 \[\frac{(a^4-1)(a^4+1)}{(a^4 + 1)(a^2+1)(a+1)}\]
 \[\frac{a^4-1}{(a^2+1)(a+1)}\]
 \[\frac{(a^2+1)(a^2-1)}{(a^2+1)(a+1)}\]
 \[\frac{a^2-1}{a+1}\]
 \[\frac{(a-1)(a+1)}{a+1}\]
 \[a-1\].

 Therefore, $x = \sqrt[16]{5} - 1$, and $(x+1)^{48} = (\sqrt[16]{5})^{48} = 5^3 = \boxed{125}$.

 





\textbf{Solution 3 (Bashing)}

Let $y=\sqrt[16]{5}$, then expanding the denominator results in:
 \[(y^8+1)(y^4+1) =(y^{12}+y^8+y^4+1)\]
 \[(y^{12}+y^8+y^4+1)(y^2+1)=(y^{14}+y^{12}+y^{10}+y^8+y^6+y^4+y^2+1)\]
 \[(y^{14}+y^{12}+y^{10}+y^8+y^6+y^4+y^2+1)(y+1) = (y^{15}+y^{14}+y^{13}+y^{12}+y^{11}+y^{10}+y^9+y^8+y^7+y^6+y^5+y^4+y^3+y^2+y+1)=\frac{y^{16}-1}{y-1}\]

 Therefore:
 \[\frac{4}{(y^{16}-1)/(y-1)} = \frac{4(y-1)}{y^{16}-1}=\frac{4(y-1)}{4} = y-1\]

 It is evident that $x+1 = (y-1)+1 = \sqrt[16]5$  as Solution 1 states.

 





\textbf{Solution 4}

Like Solution $3$, let $z=\sqrt[16]{5}$Then, the expression becomes

 $x=\frac{4}{(z+1)(z^2+1)(z^4+1)(z^8+1)}$Now, multiplying by the conjugate of each binomial in the denominator, we obtain...

 $x=\frac{4(z-1)(z^2-1)(z^4-1)(z^8-1)}{(z^2-1)(z^4-1)(z^8-1)(z^{16}-1)}\implies x=\frac{4(z-1)}{z^{16}-1}$Plugging back in,

 $x=\frac{4({\sqrt[16]{5}-1)}}{5-1}\implies x=\sqrt[16]{5}-1$

 Hence, after some basic exponent rules, we find the answer is $\boxed{125}$

 




	\item (\textit{AIME II 2005}) Circles $C_1$ and $C_2$ are externally tangent, and they are both internally tangent to circle $C_3.$ The radii of $C_1$  and $C_2$ are 4 and 10, respectively, and the centers of the three circles are all collinear. A chord of $C_3$ is also a common external tangent of $C_1$ and $C_2.$ Given that the length of the chord is $\frac{m\sqrt{n}}p$ where $m,n,$ and $p$ are positive integers, $m$ and $p$ are relatively prime, and $n$ is not divisible by the square of any prime, find $m+n+p.$


 
 \textbf{Answer: }405



\textbf{Solution}


                    \begin{figure}[H]
                        \centering
                        \includegraphics[width=0.5\textwidth]{Solution/AIME_2005_II_8_Solution_Figure_1.png}
                    \end{figure}
                    ————Let $O_1, O_2, O_3$ be the centers and $r_1 = 4, r_2 = 10,r_3 = 14$ the radii of the circles $C_1, C_2, C_3$. Let $T_1, T_2$ be the points of tangency from the common external tangent of $C_1, C_2$, respectively, and let the extension of $\overline{T_1T_2}$ intersect the extension of $\overline{O_1O_2}$ at a point $H$. Let the endpoints of the chord/tangent be $A,B$, and the foot of the perpendicular from $O_3$ to $\overline{AB}$ be $T$. From the similar \href{https://artofproblemsolving.com/wiki/index.php?title=Right\_triangle}{right triangles} $\triangle HO_1T_1 \sim \triangle HO_2T_2 \sim \triangle HO_3T$,

 
 \[\frac{HO_1}{4} = \frac{HO_1+14}{10} = \frac{HO_1+10}{O_3T}.\]

 It follows that $HO_1 = \frac{28}{3}$, and that $O_3T = \frac{58}{7}\dagger$. By the \href{https://artofproblemsolving.com/wiki/index.php?title=Pythagorean\_Theorem}{Pythagorean Theorem} on $\triangle ATO_3$, we find that 

 
 \[AB = 2AT = 2\left(\sqrt{r_3^2 - O_3T^2}\right) = 2\sqrt{14^2 - \frac{58^2}{7^2}} = \frac{8\sqrt{390}}{7}\]

 and the answer is $m+n+p=\boxed{405}$.

 $\dagger$ Alternatively, drop an altitude from $O_1$ to $O_3T$ at $O_3'$, and to $O_2T_2$ at $O_2'$. Then, $O_2O_2'=10-4=6$, and $O_1O_2=14$. But $O_1O_3O_3'$ is similar to $O_1O_2O_2'$ so $O_3O_3'=\frac{6}{14} \cdot 10=\frac{30}{7}$. From rectangles, $O_3'T=O_1T_1=4$ so $O_3T=4+\frac{30}{7}=\frac{58}{7}$.

 





\textbf{Solution 2}

Call our desired length $x$. Note for any $X$ on $\overline{AB}$ and $Y$ on $\overline{O_1O_2}$ such that $\overline{XY}\perp\overline{AB}$ that the function $f$ such that $f(\overline{O_1Y})=\overline{XY}$ is linear. Since $(0,4)$ and $(14,10)$, we can quickly interpolate that $f(10)=\overline{O_3T}=\frac{58}{7}$. Then, extend $\overline{O_3T}$ until it reaches the circle on both sides; call them $P,Q$. By Power of a Point,$\overline{PT}\cdot\overline{TQ}=\overline{AT}\cdot\overline{TB}$.Since $\overline{AT}=\overline{TB}=\frac{1}{2}x$,
 \[(\overline{PO_3}-\overline{O_3T})(\overline{QO_3}+\overline{O_3T})=\frac{1}{4}x^2\]
 \[\left(14+\frac{58}{7}\right)\left(14-\frac{58}{7}\right)=\frac{1}{4}x^2\]After solving for $x$, we get $x=\frac{8\sqrt{390}}{7}$, so our answer is $8+390+7=\boxed{405}$

 \textbf{Note} this is similar to another aime problem: \href{https://artofproblemsolving.com/wiki/index.php?title=1995\_AIME\_Problems/Problem\_4}{1995 AIME Problems/Problem 4}. This one just requires more work than the other, but the solution is practically identical.

 \textbf{Note to Ambitious Learners} Yes, we have found the length of the common external tangent. What is the length of the chord that is the common internal tangent of $C_1$ and $C_2$?

 




	\item (\textit{AIME II 2005}) For how many positive integers $n$ less than or equal to 1000 is $(\sin t + i \cos t)^n = \sin nt + i \cos nt$ true for all real $t$?


 
 \textbf{Answer: }250



\textbf{Solution 1}

We know by \href{https://artofproblemsolving.com/wiki/index.php?title=De\_Moivre\%27s\_Theorem}{De Moivre's Theorem} that $(\cos t + i \sin t)^n = \cos nt + i \sin nt$ for all \href{https://artofproblemsolving.com/wiki/index.php?title=Real\_number}{real numbers} $t$ and all \href{https://artofproblemsolving.com/wiki/index.php?title=Integer}{integers} $n$.  So, we'd like to somehow convert our given expression into a form from which we can apply De Moivre's Theorem.  

 Recall the \href{https://artofproblemsolving.com/wiki/index.php?title=Trigonometric\_identities}{trigonometric identities} $\cos \left(\frac{\pi}2 - u\right) = \sin u$ and $\sin \left(\frac{\pi}2 - u\right) = \cos u$ hold for all real $u$.  If our original equation holds for all $t$, it must certainly hold for $t = \frac{\pi}2 - u$.  Thus, the question is equivalent to asking for how many \href{https://artofproblemsolving.com/wiki/index.php?title=Positive\_integer}{positive integers} $n \leq 1000$ we have that $\left(\sin\left(\frac\pi2 - u\right) + i \cos\left(\frac\pi 2 - u\right)\right)^n = \sin n \left(\frac\pi2 -u \right) + i\cos n \left(\frac\pi2 - u\right)$ holds for all real $u$.

 $\left(\sin\left(\frac\pi2 - u\right) + i \cos\left(\frac\pi 2 - u\right)\right)^n = \left(\cos u + i \sin u\right)^n = \cos nu + i\sin nu$.  We know that two \href{https://artofproblemsolving.com/wiki/index.php?title=Complex\_number}{complex numbers} are equal if and only if both their \href{https://artofproblemsolving.com/wiki/index.php?title=Real\_part}{real part} and \href{https://artofproblemsolving.com/wiki/index.php?title=Imaginary\_part}{imaginary part} are equal.  Thus, we need to find all $n$ such that 

 




	\item (\textit{AIME II 2005}) Given that $O$ is a regular octahedron, that $C$ is the cube whose vertices are the centers of the faces of $O,$ and that the ratio of the volume of $O$ to that of $C$ is $\frac mn,$ where $m$ and $n$ are relatively prime integers, find $m+n.$


 
 \textbf{Answer: }011



\textbf{Solutions}


                    \begin{figure}[H]
                        \centering
                        \includegraphics[width=0.5\textwidth]{Solution/AIME_2005_II_10_Solution_Figure_1.png}
                    \end{figure}
                    





\textbf{Solution 1}

Let the side of the octahedron be of length $s$.  Let the \href{https://artofproblemsolving.com/wiki/index.php?title=Vertex}{vertices} of the octahedron be $A, B, C, D, E, F$ so that $A$ and $F$ are opposite each other and $AF = s\sqrt2$.  The height of the square pyramid $ABCDE$ is $\frac{AF}2 = \frac s{\sqrt2}$ and so it has volume $\frac 13 s^2 \cdot \frac s{\sqrt2} = \frac {s^3}{3\sqrt2}$ and the whole octahedron has volume $\frac {s^3\sqrt2}3$.  

 Let $M$ be the midpoint of $BC$, $N$ be the midpoint of $DE$, $G$ be the \href{https://artofproblemsolving.com/wiki/index.php?title=Centroid}{centroid} of $\triangle ABC$ and $H$ be the centroid of $\triangle ADE$.  Then $\triangle AMN \sim \triangle AGH$ and the symmetry ratio is $\frac 23$ (because the \href{https://artofproblemsolving.com/wiki/index.php?title=Triangle\_median}{medians} of a triangle are trisected by the centroid), so $GH = \frac{2}{3}MN = \frac{2s}3$.  $GH$ is also a diagonal of the cube, so the cube has side-length $\frac{s\sqrt2}3$ and volume $\frac{2s^3\sqrt2}{27}$.  The ratio of the volumes is then $\frac{\left(\frac{s^3\sqrt2}{3}\right)}{\left(\frac{2s^3\sqrt2}{27}\right)} = \frac92$ and so the answer is $\boxed{011}$.

 





\textbf{Solution 2}

Let the octahedron have vertices $(\pm 3, 0, 0), (0, \pm 3, 0), (0, 0, \pm 3)$.  Then the vertices of the cube lie at the centroids of the faces, which have coordinates $(\pm 1, \pm 1, \pm 1)$.  The cube has volume 8.  The region of the octahedron lying in each octant is a \href{https://artofproblemsolving.com/wiki/index.php?title=Tetrahedron}{tetrahedron} with three edges mutually perpendicular and of length 3.  Thus the octahedron has volume $8 \cdot \left(\frac 16 \cdot3^3\right) = 36$, so the ratio is $\frac {36}{8} = \frac 92$ and so the answer is $\boxed{011}$.

 





\textbf{Solution 3}

Continue Solution 2 until realizing the volume of the cube is 8. The side length of the octahedron through distance formulae is $3\sqrt{2}$. One can quickly derive the volume of a octahedron by multiplying two times the volume of a tetrahedron. The height of the top tetrahedron is 3, and the area of the square is 18, so the volume of the top tetrahedron is 18, and the total volume of the octahedron is 36. So, the ratio is $\frac {36}{8} = \frac{9}{2}$ and so the answer is $\boxed{011}$

 




	\item (\textit{AIME II 2005}) Let $m$ be a positive integer, and let $a_0, a_1,\ldots,a_m$ be a sequence of reals such that $a_0 = 37, a_1 = 72, a_m = 0,$ and $a_{k+1} = a_{k-1} - \frac 3{a_k}$ for $k = 1,2,\ldots, m-1.$ Find $m.$


 
 \textbf{Answer: }889



\textbf{Solution 1}

For $0 < k < m$, we have

 <center>$a_{k}a_{k+1} = a_{k-1}a_{k} - 3$.

 </center>Thus the product $a_{k}a_{k+1}$ is a \href{https://artofproblemsolving.com/wiki/index.php?title=Monovariant}{monovariant}: it decreases by 3 each time $k$ increases by 1.  For $k = 0$ we have $a_{k}a_{k+1} = 37\cdot 72$, so when $k = \frac{37 \cdot 72}{3} = 888$, $a_{k}a_{k+1}$ will be zero for the first time, which implies that $m = \boxed{889}$, our answer.

 
Note: In order for $a_{m} = 0$ we need $a_{m-2}a_{m-1}=3$ simply by the recursion definition.

 





\textbf{Solution 2}

Plugging in $k = m-1$ to the given relation, we get $0 = a_{m-2} - \frac {3}{a_{m-1}} \implies{a_{m-2}a_{m-1} = 3}$. Inspecting the value of $a_{k}a_{k+1}$ for small values of $k$, we see that $a_{k}a_{k+1} = 37\cdot 72 - 3k$. Setting the RHS of this equation equal to $3$, we find that $m$ must be $\boxed{889}$.

 \textbf{Induction Proof} As above, we experiment with some values of $a_{k}a_{k+1}$, conjecturing that $a_{m-p}a_{m-p-1}$ = $3p$ ,where $m$ is a positive integer and so is $p$, and we prove this formally using induction. The base case is for $p = 1$, $a_{m} = a_{m-2} - 3/a_{m-1}$ Since $a_{m} = 0$, $a_{m-1}a_{m-2} = 3$; from the recursion given in the problem $a_{m-p+1} = a_{m-p-1} - 3/a_{m-p}$, so $a_{m-p+1} = 3p/a_{m-p} - 3/a_{m-p} = 3(p-1)/a_{m-p}$, so $a_{m-p}a_{m-p+1} = a_{m-(p-1)}a_{m-(p-1)-1} = 3(p-1)$, hence proving our formula by induction.~USAMO2023

 





\textbf{Solution 3 (Telescoping)}

Note that $a_{k+1}\cdot a_{k} - a_{k}\cdot a_{k-1} = - 3$. Then, we can generate a sum of series of equations such that $\sum_{k=1}^{m-1} a_{k+1}\cdot a_{k} - a_{k}\cdot a_{k-1} = - 3(m-1)$. Then, note that all but the first and last terms on the LHS cancel out, leaving us with $a_m\cdot a_{m-1} - a_1\cdot a_0 = -3(m-1)$. Plugging in $a_m=0$, $a_0=37$, $a_1=72$, we have $-37\cdot 72 = -3(m-1)$. Solving for $m$ gives $m=\boxed{889}$. ~sigma

 





\textbf{Solution 4 (Explicit, Elementary Mathematics)}


 \[a_{k+1} = a_{k-1} - \frac{3}{a_k}\]
 \[a_{k+1}a_{k} = a_{k-1}a_{k} - 3\]

 Define $b_n = a_{k}a_{k+1}$. Then,

 
 \[b_n = b_{n-1} - 3\].

 Note that $b_0 = 2664$, so the explicit relation for $b$ is $b_n = 2664-3n$. 

 We don't see a clean homogenous equation, so we end our recurrence seeking here. However, now we have enough information. Because $a_n(a_{n+1}) = 2664-3n$, we just need the $m$ such that $a_m(a_{m+1}) = 2664-3m$. However, $a_{m+1}$ is undefined. Therefore, we must consider $a_{m-1}a_{m} = 2664-3(m-1)$. We know $a_m = 0$, so $0 = 2664-3(m-1)$, implying $m = \boxed{889}$.

 





\textbf{Solution 5 (Explicit, Olympiad Mathematics)}

$\textbf{Lemma}:$

 
 \[a_{2n} = 37 \prod_{i=0}^{n-1} \frac{887-2i}{888 - 2i}\]
 \[a_{2n+1} = 72 \prod_{i=0}^{n-1} \frac{886-2i}{887-2i}\].
 \[a_m = 0 \ \text{if and only if} \ m = \boxed{889}\]

 $\textbf{Proof of the Lemma}:$

 
 \[a_{k+1} = a_{k-1} - \frac{3}{a_k}\]
 \[a_{k+1}a_{k} = a_{k-1}a_{k} - 3\]

 Define $b_n = a_{k}a_{k+1}$. Then,

 
 \[b_n = b_{n-1} - 3\].

 Note that $b_0 = 2664$, so the explicit relation for $b$ is $b_n = 2664-3n$. Now, define $a_{n+1} = \frac{b_n}{a_n}$. The value of $a_n$ will change one whether or not the index is even or odd. We see a cascading cancellation, so we have

 
 \[a_{2n} = a_0 \cdot \frac{b_1}{b_0} \cdot \frac{b_3}{b_2} \cdot \dots \cdot \frac{b_{2n-1}}{b_{2n-2}} = 37 \prod_{i=0}^{n-1} \frac{887-2i}{888 - 2i}\]

 
 \[a_{2n+1} = a_1 \cdot \frac{b_2}{b_1} \cdot \frac{b_4}{b_3} \cdot \dots \cdot \frac{b_{2n}}{b_{2n-1}} = 72 \prod_{i=0}^{n-1} \frac{886-2i}{887-2i}\].

 We require the first term that makes the product 0, which will only happen if one of the terms in the sequence within the product simplifies to 0. In that case, it is required that the requested $a_m$ term has $m$ as an odd index, as otherwise no such product can exist (as then $m$ is not an integer). So, we must have that $m$ is an odd index, and that $886-2i = 0 \implies i = 443 \ \text{where} \ i = n-1$. Here, we have that the unique $n$ is $n = 444$, so the requested $m$ is then $m = 2(444) + 1 = \boxed{889}$.

 




	\item (\textit{AIME II 2005}) Square $ABCD$ has center $O, AB=900, E$ and $F$ are on $AB$ with $AE<BF$ and $E$ between $A$ and $F, m\angle EOF =45^\circ,$ and $EF=400.$ Given that $BF=p+q\sqrt{r},$ where $p,q,$ and $r$ are positive integers and $r$ is not divisible by the square of any prime, find $p+q+r.$


 
 \textbf{Answer: }307



\textbf{Solution 1 (trigonometry)}


                    \begin{figure}[H]
                        \centering
                        \includegraphics[width=0.5\textwidth]{Solution/AIME_2005_II_12_Solution_Figure_1.png}
                    \end{figure}
                    Let $G$ be the foot of the \href{https://artofproblemsolving.com/wiki/index.php?title=Perpendicular}{perpendicular} from $O$ to $AB$. Denote $x = EG$ and $y = FG$, and $x > y$ (since $AE < BF$ and $AG = BG$). Then $\tan \angle EOG = \frac{x}{450}$, and $\tan \angle FOG = \frac{y}{450}$.

 By the \href{https://artofproblemsolving.com/wiki/index.php?title=Trigonometric\_identity}{tangent addition rule} $\left( \tan (a + b) = \frac{\tan a + \tan b}{1 - \tan a \tan b} \right)$, we see that 
 \[\tan 45 = \tan (EOG + FOG) = \frac{\frac{x}{450} + \frac{y}{450}}{1 - \frac{x}{450} \cdot \frac{y}{450}}.\] Since $\tan 45 = 1$, this simplifies to $1 - \frac{xy}{450^2} = \frac{x + y}{450}$. We know that $x + y = 400$, so we can substitute this to find that $1 - \frac{xy}{450^2} = \frac 89 \Longrightarrow xy = 150^2$.

 Substituting $x = 400 - y$ again, we know have $xy = (400 - y)y = 150^2$. This is a quadratic with roots $200 \pm 50\sqrt{7}$. Since $y < x$, use the smaller root, $200 - 50\sqrt{7}$.

 Now, $BF = BG - FG = 450 - (200 - 50\sqrt{7}) = 250 + 50\sqrt{7}$. The answer is $250 + 50 + 7 = \boxed{307}$.

 





\textbf{Solution 2 (synthetic)}


                    \begin{figure}[H]
                        \centering
                        \includegraphics[width=0.5\textwidth]{Solution/AIME_2005_II_12_Solution_Figure_1.png}
                    \end{figure}
                    Label $BF=x$, so $EA =$ $500 - x$. Rotate $\triangle{OEF}$ about $O$ until $EF$ lies on $BC$. Now we know that $\angle{EOF}=45^\circ$ therefore $\angle BOF+\angle AOE=45^\circ$ also since $O$ is the center of the square.  Label the new triangle that we created $\triangle OGJ$.  Now we know that rotation preserves angles and side lengths, so $BG=500-x$ and $JC=x$.  Draw $GF$ and $OB$. Notice that $\angle BOG =\angle OAE$ since rotations preserve the same angles so $\angle{FOG}=45^\circ$ too. By SAS we know that $\triangle FOE\cong \triangle FOG,$ so $FG=400$. Now we have a right $\triangle BFG$ with legs $x$ and $500-x$ and hypotenuse $400$. By the \href{https://artofproblemsolving.com/wiki/index.php?title=Pythagorean\_Theorem}{Pythagorean Theorem}, 

 
 \begin{align*} (500-x)^2+x^2&=400^2 \\ 250000-1000x+2x^2&=16000 \\ 90000-1000x+2x^2&=0 \end{align*}


 and applying the \href{https://artofproblemsolving.com/wiki/index.php?title=Quadratic\_formula}{quadratic formula} we get that $x=250\pm 50\sqrt{7}$.  Since $BF > AE,$ we take the positive root, and our answer is $p+q+r = 250 + 50 + 7 = 307$.

 





\textbf{Solution 3 (similar triangles)}


                    \begin{figure}[H]
                        \centering
                        \includegraphics[width=0.5\textwidth]{Solution/AIME_2005_II_12_Solution_Figure_1.png}
                    \end{figure}
                    Let the midpoint of $\overline{AB}$ be $M$ and let $FB = x$, so then $MF = 450 - x$ and $AF = 900 - x$. Drawing $\overline{AO}$, we have $\triangle OEF\sim\triangle AOF$, so
 \[\frac{OF}{EF} = \frac{AF}{OF}\Rightarrow (OF)^2 = 400(900 - x).\]By the Pythagorean Theorem on $\triangle OMF$,
 \[(OF)^2 = 450^2 + (450 - x)^2.\]Setting these two expressions for $(OF)^2$ equal and solving for $x$ (it is helpful to scale the problem down by a factor of 50 first), we get $x = 250\pm 50\sqrt{7}$. Since $BF > AE$, we want the value $x = 250 + 50\sqrt{7}$, and the answer is $250 + 50 + 7 = \boxed{307}$.

 





\textbf{Solution 4 (Abusing Stewart)}

Let $x = BF$, so $AE = 500-x$. Let $a = OE$, $b = OF$. Applying Stewart's Theorem on triangles $AOB$ twice, first using $E$ as the base point and then $F$, we arrive at the equations 
 \[(450 \sqrt{2})^2 (900) = 900(500-x)(400+x) + a^2 (900)\] and 
 \[(450 \sqrt{2})^2 (900) = 900x(900-x) + b^2 (900)\] Now applying law of sines and law of cosines on $\triangle EOF$ yields 
 \[\frac{1}{2} ab \sin 45^{\circ}  = \frac{4}{9} \times \frac{1}{4}  \times 900^2 = 202500\] and 
 \[a^2+b^2- 2 ab \cos 45^{\circ} = 160000\] Solving for $ab$ from the sines equation and plugging into the law of cosines equation yields $a^2+b^2 = 290000$. We now finish by adding the two original stewart equations and obtaining: 
 \[2(450\sqrt{2})^2 = (500-x)(400+x)+x(900-x)+520000\] This is a quadratic which only takes some patience to solve for $x = 250 + 50\sqrt{7}$

 





\textbf{Solution 5 (Complex Numbers)}

Let lower case letters be the complex numbers correspond to their respective upper case points in the complex plane, with $o = 0, a = -450 + 450i, b = 450 + 450i$, and $f = x + 450i$. Since $EF$ = 400, $e = (x-400) + 450i$. From $\angle{EOF} = 45^{\circ}$, we can deduce that the rotation of point $F$ 45 degrees counterclockwise, $E$, and the origin are collinear. In other words, 
 \[\dfrac{e^{i \frac{\pi}{4}} \cdot (x + 450i)}{(x - 400) + 450i}\] is a real number. Simplyfying using the fact that $e^{i \frac{\pi}{4}} = \dfrac{\sqrt{2}}{2} + i \dfrac{\sqrt{2}}{2}$, clearing the denominator, and setting the imaginary part equal to $0$, we eventually get the quadratic 
 \[x^2 - 400x + 22500 = 0\] which has solutions $x = 200 \pm 50\sqrt{7}$. It is given that $AE < BF$, so $x = 200 - 50\sqrt{7}$ and 
 \[BF = 450 - (200 - 50\sqrt{7}) = 250 + 50\sqrt{7} \Rightarrow \boxed{307}.\]

 -MP8148

 





\textbf{Solution 6}


                    \begin{figure}[H]
                        \centering
                        \includegraphics[width=0.5\textwidth]{Solution/AIME_2005_II_12_Solution_Figure_1.png}
                    \end{figure}
                    Let G be a point such that it lies on AB, and GOE is 90 degrees. Let H be foot of the altitude from O to AB.Since $\triangle GOE \sim \triangle OHE$, $\frac{GO}{OE} = \frac{450}{x}$, and by \href{https://artofproblemsolving.com/wiki/index.php?title=Angle\_Bisector\_Theorem}{Angle Bisector Theorem}, $\frac{GF}{FE} = \frac{450}{x}$. Thus, $GF = \frac{450 \cdot 400}{x}$. $AF = AH-FH = 50+x$, and $KA = EB$ (90 degree rotation), and now we can bash on 2 similar triangles $\triangle GAK \sim \triangle GHO$.

 
 \[\frac{GA}{AK} = \frac{GH}{OH}\] 
 \[\frac{\frac{450 \cdot 400}{x}-50-x}{450-x} = \frac{\frac{450 \cdot 400}{x}+400-x}{450}\]I hope you like expanding
 \[x^2 - 850x + \frac{81000000}{x} = -450x - 22500 + \frac{81000000}{x}\]
 \[x^2 - 400x + 22500 = 0\]Quadratic formula gives us
 \[x = 200 \pm 50 \sqrt{7}\]Since AE < BF 
 \[x = 200 - 50 \sqrt{7}\]Thus, 
 \[BF = 250 + 50 \sqrt{7}\]So, our answer is $\boxed{307}$.

 -AlexLikeMath

 





\textbf{Solution 7 (Using a Circle)}


                    \begin{figure}[H]
                        \centering
                        \includegraphics[width=0.5\textwidth]{Solution/AIME_2005_II_12_Solution_Figure_1.png}
                    \end{figure}
                    We know that G is on the perpendicular bisector of $EF$, which means that $EJ=JF=200$, $EG=GF=200\sqrt{2}$ and $GH=250$. Now, let $HO$ be equal to $x$. We can set up an equation with the Pythagorean Theorem: 

 
 \begin{align*} \sqrt{x^2+250^2}&=(200\sqrt{2})^2 \\ x^2+62500&=80000 \\ x^2&=17500 \\ x&=50\sqrt{7} \end{align*}


 Now, since $IO=450$, 

 
 \begin{align*} HI&=450-x \\ &=450-50\sqrt{7} \\ \end{align*}
 \\

 Since $HI=AJ$, we now have:

 
 \begin{align*} BF&=AB-AJ-JF \\ &=900-(450-50\sqrt{7})-200 \\ &=250+50\sqrt{7} \\ \end{align*}
 

 This means that our answer would be $250+50+7=\boxed{307}$

 





\textbf{Solution 8 (More Similar Triangles)}

Construct $BO, AO.$ Let $\angle{FOB} = \alpha.$ Also let $FB = x$ then $AE = 500-x.$ We then have from simple angle-chasing:
 \begin{align*} \angle{BFO} = 135 - \alpha \\ \angle{OFE} = 45 + \alpha \\ \angle{EOA} = 45 - \alpha \\ \angle{AEO} = 90 + \alpha \\ \angle{OEF} = 90 - \alpha. \end{align*}
From AA similarity we have 
 \[\triangle{EOB} \sim \triangle{EFO}.\] This gives the ratios, 
 \[\dfrac{400 + x}{EO} = \dfrac{450\sqrt{2}}{FO}.\]Similarly from AA similarity 
 \[\triangle{FOA} \sim \triangle{FEO}.\] So we get the ratios 
 \[\dfrac{EO}{450\sqrt{2}} = \dfrac{FO}{900-x}.\] We can multiply to get 
 \[\dfrac{400 + x}{450\sqrt{2}} = \dfrac{450\sqrt{2}}{900 - x}.\] Cross-multiplying reveals
 \[360000 + 500x - x^2 = 405000.\] Bringing everything to one side we have 
 \[x^2 - 500x + 45000 = 0.\] By the quadratic formula we get 
 \[x = \dfrac{500 + \sqrt{500^2 - 4\cdot45000}}{2} = \dfrac{500 + \sqrt{70000}}{2} = \dfrac{500 + 100\sqrt{7}}{2} = 250 + 50\sqrt{7}.\]Therefore 
 \[p+q+r = 250 + 50 + 7 = \boxed{307}.\]~aa1024

 


 





\textbf{Solution 9}

We use ratio lemma and Stewart's theorem:Connect $OA, OE, OF, OB$ and let $AE = x$ and $BF = 500 - x.$Let angle $AOE = y,$ hence $BOF = 45 - y.$Now, we apply Stewart's theorem in triangles $AOF$ and $BOE$ to get $OE$ and $OF$ in terms of $x$finally, calculate $x/400$ and $500-x/400$ using ratio lemma to find $x$ and $y$

 





\textbf{Solution 10(Similar Triangles)}

Draw AO, OB, and extend OB to D. Let $\angle{FOB} = \alpha.$ Then, after angle chasing, we find that 
 \[\angle{AEB} = 90 + \alpha\]. Using this, we draw a line perpendicular to $AB$ at $E$ to meet $BD$ at $M$. Since $\angle{MEO} = \alpha$ and $\angle{EMO} = 45$, we have that 
 \[\triangle{EMO} \sim \triangle{OBF}\]Let $FB = x$. Then $EM = 400+x$. Since $FB/BO = \frac{x}{450\sqrt{2}}$, and $MO/EM = FB/OB$, we have 
 \[MO = \frac{(400+x)x}{450\sqrt{2}}\] Since $\triangle{EBM}$ is a $45-45-90$ triangle, 
 \[(400+x)\sqrt{2} = 450 \sqrt{2} + \frac{(400+x)x}{450\sqrt{2}}\]Solving for $x$, we get that $x=250 +- 50s\sqrt{7}$, but since $FB>AE$, $FB = 250+50\sqrt{7}$, thus 
 \[p+q+r=\boxed{307}\]-dchang0524

 




	\item (\textit{AIME II 2005}) Let $P(x)$ be a polynomial with integer coefficients that satisfies $P(17)=10$ and $P(24)=17.$ Given that $P(n)=n+3$ has two distinct integer solutions $n_1$ and $n_2,$ find the product $n_1\cdot n_2.$


 
 \textbf{Answer: }418



\textbf{Solution}

We define $Q(x)=P(x)-x+7$, noting that it has roots at $17$ and $24$. Hence $P(x)-x+7=A(x-17)(x-24)$. In particular, this means that $P(x)-x-3=A(x-17)(x-24)-10$. Therefore, $x=n_1,n_2$ satisfy $A(x-17)(x-24)=10$, where $A$, $(x-17)$, and $(x-24)$ are integers. This cannot occur if $x\le 17$ or $x\ge 24$ because the product $(x-17)(x-24)$ will either be too large or not be a divisor of $10$. We find that $x=19$ and $x=22$ are the only values that allow $(x-17)(x-24)$ to be a factor of $10$.  Hence the answer is $19\cdot 22=\boxed{418}$.

 


 





\textbf{Solution 2}

We know that $P(n)-(n+3)=0$ so $P(n)$ has two distinct solutions so $P(x)$ is at least quadratic. Let us first try this problem out as if $P(x)$ is a quadratic polynomial. Thus $P(n)-(n+3)= an^2+(b-1)n+(c-3)=0$ because $P(n)=an^2+bn+c$ where $a,b,c$ are all integers. Thus $P(x)=ax^2+bx+c$ where $a,b,c$ are all integers. We know that $P(17)$ or $289a+17b+c=10$ and $P(24)$ or $576a+24b+c=17$. By doing $P(24)-P(17)$ we obtain that $287a+7b=7$ or $41a+b=1$ or $-41a=b-1$. Thus $P(n)=an^2- (41a)n+(c-3)=0$.  Now we know that $b=-41a+1$, we have $289a+17(-41a+1)+c=10$ or $408a=7+c$ which makes $408a-10=c-3$. Thus $P(n)=an^2-(41a)n+(408a-10)=0$. By Vieta's formulas, we know that the sum of the roots($n$) is equal to 41 and the product of the roots($n$) is equal to $408-\frac{10}{a}$. Because the roots are integers $\frac{10}{a}$ has to be an integer, so $a=1,2,5,10,-1,-2,-5,-10$. Thus the product of the roots is equal to one of the following: $398,403,406,407,409,410,413,418$. Testing every potential product of the roots, we find out that the only product that can have divisors that sum up to $41$ is $\boxed{418}$.

 -David Camacho

 





\textbf{Solution 3}

We have $P(n_1) = n_1+3$. Using the property that $a - b \mid P(a) - P(b)$ whenever $a$ and $b$ are distinct integers, we get 
 \[n_1 - 17 \mid P(n_1) - P(17) = (n_1+3) - 10 = n_1 - 7,\]and 
 \[n_1 - 24 \mid P(n_1) - P(24) = (n_1+3)-17=n_1-14.\]Since $n_1 - 7 = 10 + (n_1-17)$ and $n_1-14 = 10 + (n_1-24)$, we must have 
 \[n_1 - 17 \mid 10 \; \text{and} \; n_1-24 \mid 10.\]We look for two divisors of $10$ that differ by $7$; we find that $\{2, -5\}$ and $\{5, -2\}$ satisfies these conditions. Therefore, either $n_1 - 24 = -5$, giving $n_1 = 19$, or $n_1 - 24 = -2$, giving $n_1 = 22$. From this, we conclude that $n_1, n_2 = \boxed{19, 22}$ and that $n_1\cdot n_2=\boxed{418}$.

 





\textbf{Solution 4}

Based on the second part of the problem statement, $P(n)$ is at least a quadratic. So, we can test that out. Plugging in our known values, we know that $P(24)=576a+24b+c=17$ and $P(17)=289a+17b+c=10$. So, $P(24)-P(17)=287a+7b=7$ and $41a+b=1$. Plugging in the working $n_1$ and $n_2$ values, we get that $an_1^2+bn_1+c=n_1+3$ and $an_2^2+bn_2+c=n_2+3$. Subtracting the two equations gives that $a(n_1^2-n_2^2)+b(n_1-n_2)=n_1-n_2 \implies a(n_1+n_2)+b=1=41a+b \implies n_1+n_2=41$.  Adding the two previous equations gives that $47=2c+41b+a(n_1^2+n_2^2)$ and adding the two equations we were given gives us $27=2c+41b+865a$. Subtracting implies that $20=a(n_1^2+n_2^2-865)=a((n_1+n_2)^2-2n_1 \cdot n_2-865) = a(1681-865-2n_1 \cdot n_2)=a(816-2n_1 \cdot n_2)$. 

 So, since both parts of the expression have to be integers, $816-2n_1 \cdot n_2=2,4,10,20,-2,-4,-10,-20 \implies n_1 \cdot n_2 = 407,406,403,398,409,410,413,418$. 

 Plugging in values that satisfy $n_1+n_2=41$ gives us that the only possible solution is $n_1\cdot n_2=\boxed{418}$.

 




	\item (\textit{AIME II 2005}) In triangle $ABC, AB=13, BC=15,$ and $CA = 14.$ Point $D$ is on $\overline{BC}$ with $CD=6.$ Point $E$ is on $\overline{BC}$ such that $\angle BAE\cong \angle CAD.$ Given that $BE=\frac pq$ where $p$ and $q$ are relatively prime positive integers, find $q.$


 
 \textbf{Answer: }463



\textbf{Solution 1}


                    \begin{figure}[H]
                        \centering
                        \includegraphics[width=0.5\textwidth]{Solution/AIME_2005_II_14_Solution_Figure_1.png}
                    \end{figure}
                    By the \href{https://artofproblemsolving.com/wiki/index.php?title=Law\_of\_Sines}{Law of Sines} and since $\angle BAE = \angle CAD, \angle BAD = \angle CAE$, we have

 
 \begin{align*}  \frac{CD \cdot CE}{AC^2} &= \frac{\sin CAD}{\sin ADC} \cdot \frac{\sin CAE}{\sin AEC}  \\ &= \frac{\sin BAE \sin BAD}{\sin ADB \sin AEB} \\ &= \frac{\sin BAE}{\sin AEB} \cdot \frac{\sin BAD}{\sin ADB}\\ &= \frac{BE \cdot BD}{AB^2} \end{align*}


 Substituting our knowns, we have $\frac{CE}{BE} = \frac{3 \cdot 14^2}{2 \cdot 13^2} = \frac{BC - BE}{BE} = \frac{15}{BE} - 1 \Longrightarrow BE = \frac{13^2 \cdot 15}{463}$. The answer is $q = \boxed{463}$.

 





\textbf{Solution 2 (Similar Triangles)}

Drop the altitude from A and call the base of the altitude Q. Also, drop the altitudes from E and D to AB and AC respectively. Call the feet of the altitudes R and S respectively. 

 From here, we can use Heron's Formula to find the altitude. The area of the triangle is $\sqrt{21*6*7*8} = 84$. We can then use similar triangles with triangle $AQC$ and triangle $DSC$ to find $DS=\frac{24}{5}$. Consequently, from Pythagorean theorem, $SC = \frac{18}{5}$ and $AS = 14-SC = \frac{52}{5}$. We can also use the Pythagorean theorem on triangle $AQB$ to determine that $BQ = \frac{33}{5}$.

 Label $AR$ as $y$ and $RE$ as $x$. $RB$ then equals $13-y$. Then, we have two similar triangles. 

 Firstly: $\triangle ARE \sim \triangle ASD$. From there, we have $\frac{x}{y} = \frac{\frac{24}{5}}{\frac{53}{5}}$. 

 Next: $\triangle BRE \sim \triangle BQA$. From there, we have $\frac{x}{13-y} = \frac{\frac{56}{5}}{\frac{33}{5}}$.

 Solve the system to get $x = \frac{2184}{463}$ and $y = \frac{4732}{463}$.  Notice that 463 is prime, so even though we use the Pythagorean theorem on $x$ and $13-y$, the denominator won't change. The answer we desire is $\boxed{463}$.

 





\textbf{Solution 3 (LoC and LoS bash)}

Let $\angle CAD = \angle BAE = \theta$. Note by Law of Sines on $\triangle BEA$ we have
 \[\frac{BE}{\sin{\theta}} = \frac{AE}{\sin{B}} = \frac{AB}{\sin{\angle BEA}}\]As a result, our goal is to find $\sin{\angle BEA}$ and $\sin{\theta}$ (we already know $AB$). 

 Let the foot of the altitude from $A$ to $BC$ be $H$. By law of cosines on $\triangle ABC$ we have
 \[169 = 196 + 225 - 2 \cdot 14 \cdot 15 \cdot \cos{C} \Rightarrow \cos{C} = \frac{3}{5}\]It follows that $AH = \frac{56}{5}$ and $HC = \frac{42}{5} \Rightarrow HD = \frac{12}{5}$. 

 Note that by PT on $\triangle AHD$ we have that $AD^2 = \left(\frac{56}{5}\right)^2 + \left(\frac{12}{5}\right)^2 = \frac{656}{5}$. By Law of Sines on $\triangle ADC$ (where we square everything to avoid taking the square root) we see
 \[\frac{36}{\sin^2{\theta}} = \frac{656}{5 \cdot \frac{16}{25}} \Rightarrow \sin^2{\theta} = \frac{36}{205}.\]How are we going to find $\sin{\angle BEA}$ though? $\angle BEA$ and $\theta$ are in the same triangle. Applying Law of Sines on $\triangle ABC$ we see that
 \[\frac{13}{\frac{4}{5}} = \frac{14}{\sin{\angle B}} \Rightarrow \sin{\angle B} = \frac{56}{65} \Rightarrow \cos{\angle B} = \frac{33}{65}.\]$\theta$, $\angle B$, and $\angle BEA$ are all in the same triangle. We know they add up to $180^{\circ}$. There's a good chance we can exploit this using the identity $\sin{p} = \sin{180^{\circ}-p}$.

 We have that $\sin{(180^{\circ} - (\theta + \angle B))} = \sin{\angle BEA} = \sin{(\theta + \angle B)}$. Success! We know $\sin{\theta}$ and $\sin{\angle B}$ already. Applying the $\sin$ addition formula we see
 \[\sin{\theta + \angle B} = \sin{\theta} \cos{\angle B} + \sin{\angle B} \cos{\theta} = \frac{6}{\sqrt{205}} \cdot \frac{33}{65} + \frac{56}{65} \cdot \frac{13}{\sqrt{205}}=\frac{1}{65 \cdot \sqrt{205}} (198 + 728) = \frac{926}{65 \sqrt{205}}.\]This is the last stretch! Applying Law of Sines a final time on $\triangle BEA$ we see
 \[\frac{BE}{\sin{\theta}} = \frac{13}{\sin{BEA}} \Rightarrow \frac{BE}{\frac{6}{\sqrt{205}}} = \frac{13}{\frac{926}{65\sqrt{205}}} \Rightarrow \frac{BE}{6} = \frac{13 \cdot 65}{926} \Rightarrow \frac{13 \cdot 65 \cdot 6}{926} = BE = \frac{2535}{463}.\]It follows that the answer is $\boxed{463}$.

 





\textbf{Solution 4 (Ratio Lemma and Angle Bisector Theorem)}

Let $AK$ be the angle bisector of $\angle A$ such that $K$ is on $BC$.

 Then $\angle KAB = \angle KAC$, and thus $\angle KAE = \angle KAD$.

 By the Ratio Lemma, $\frac{BE}{KE} = \frac{BA}{KA} * \frac{\sin{BAE}}{\sin{KAE}}$ and $\frac{CD}{KD} = \frac{CA}{KA} * \frac{\sin{CAD}}{\sin{KAD}}$.

 This implies that $\frac{BE}{KE*BA} = \frac{CD}{KD*CA}$.

 Thus, $\frac{BE}{KE} = \frac{13}{14} * \frac{6}{DK}$.

 $DK = CK - 6 = 14*15/27 - 6 = 16/9$. Thus, $\frac{BE}{KE} = \frac{13*54}{14*16}$.

 Additionally, $BE + KE = 13*15/27 = 65/9$. Solving gives that $q = 463.$

 Alternate:By the ratio lemma,$BD/DC = (13/14)*(\sin BAD/\sin DAC)$$EC/EB = (14/13)*(\sin EAC/\sin BAE)$

 Combining these, we get$(BD/DC)(14/13) = (EC/EB)(13/14)$$(3/2)(14/13)(14/13) = (15-x)(x)$

 $x = 2535/463$Thus, $q = 463$

 





\textbf{Solution 5 (Isogonal lines with respect to A angle bisesector)}

Since $AE$ and $AD$ are isogonal with respect to the $A$ angle bisector, we have 
 \[\frac{BE}{EC}\cdot \frac{BD}{DC}=(\frac{AB}{AC})^2.\]  To prove this, let $\angle BAE=\angle DAC=x$ and $\angle BAD=\angle CAE=y.$  Then, by the Ratio Lemma, we have 
 \[\frac{BD}{DC}=\frac{AB\sin y}{AC\sin x}\] 
 \[\frac{BE}{EC}=\frac{AB\sin x}{AC\sin y}\] and multiplying these together proves the formula for isogonal lines. Hence, we have 
 \[\frac{BE}{15-BE}\cdot \frac{9}{6}=\frac{169}{196}\implies BE=\frac{2535}{463}\] so our desired answer is $\boxed{463}.$

 


 





\textbf{Solution 6 (Tangent subtraction formulas)}

Note: We first recall some helpful tips regarding 13, 14, 15 triangles. Drawing an altitude H from B to AC results in AHB being a 5-12-13 right triangle and CHB being a 3-4-5 (9-12-15) right triangle. 

 
                    \begin{figure}[H]
                        \centering
                        \includegraphics[width=0.5\textwidth]{Solution/AIME_2005_II_14_Solution_Figure_1.png}
                    \end{figure}
                    Now we start by drawing altitudes from D and E onto AC, labeling them as F and G, and labelling $\angle DAG = \alpha$. Now we know that $\overline{DF} = \frac{24}{5}$ and $\overline{FC} = \frac{18}{5}$. Therefore, $\overline{AF} = \frac{52}{5}$, so $\tan{(\alpha)} = \frac{6}{13}$. Our goal now is to use tangent $\angle EAG$ in triangle $AEG$. We set $\overline{BE}$ to $x$, so $\overline{ED} = 9 - x$ and $\overline{EC} = 15 - x$, so $\overline{EG} = \frac{4}{5}(15-x)$ and $\overline{GC} = \frac{3}{5}(15-x)$ so $\overline{AG} = \frac{3x+25}{5}$. Now we just need tangent of $\angle EAG$. 

 We find this using $\tan{(EAG)} = \tan{(A - \alpha)} = \frac{\tan{A} - \tan{\alpha}}{1 + \tan{A}\tan{\alpha}}$, which is $\frac{\frac{12}{5} - \frac{6}{13}}{1 + \frac{12}{5} \cdot \frac{6}{13}}$ or $\frac{126}{137}$. Now we solve the equation $\tan{\angle EAG} = \frac{126}{137} = \frac{\frac{60-4x}{5}}{\frac{3x+25}{5}}$, so $x = \frac{2535}{463}$

 





\textbf{Solution 8 (Isogonal Conjugates)}

Let $ED = x$, such that $BE = 9-x$. Since $\overline{AE}$ and $\overline{AD}$ are isogonal, we get $\frac{9-x}{6+x} \cdot \frac{9}{6} = \frac{13^2}{14^2} \Rightarrow 588(9 - x) = 338(6 + x)$, and we can solve to get $x = \frac{1632}{463}$(and $BE = \frac{2535}{463}$). Hence, our answer is $\boxed{463}$. - Spacesam

 





\textbf{Solution 9 (Short and no IQ Required Altogether-Bash)}

Diagram borrowed from Solution 1.

 
                    \begin{figure}[H]
                        \centering
                        \includegraphics[width=0.5\textwidth]{Solution/AIME_2005_II_14_Solution_Figure_1.png}
                    \end{figure}
                    Applying \href{https://artofproblemsolving.com/wiki/index.php?title=Law\_of\_Cosines}{Law of Cosines} on $\bigtriangleup ABC$ with respect to $\angle C$ we have
 \[AB^2=AC^2+BC^2-2(AC)(BC)\cos C\]Solving gets $\cos C=\frac{3}{5}$, which implies that 
 \[\sin C=\sqrt{1-\cos C}=\frac{4}{5}\]Applying \href{https://artofproblemsolving.com/wiki/index.php?title=Stewart\%27s\_Theorem}{Stewart's Theorem} with \href{https://artofproblemsolving.com/wiki/index.php?title=Cevian}{cevian} $AD$ we have
 \[(BC)(AD)^2+(BC)(BD)(CD)=(CD)(AB)^2+(BD)(AC)^2\]Solving gets $AD=\frac{4\sqrt{205}}{5}$.

 Applying \href{https://artofproblemsolving.com/wiki/index.php?title=Law\_of\_Sines}{Law of Sines} on $\bigtriangleup ACD$ to solve for $\sin CAD$ we have
 \[\frac{AD}{\sin C}=\frac{CD}{\sin CAD}\]Solving gets $\sin CAD=\frac{6\sqrt{205}}{205}$. Thus $\sin BAE=\sin CAD=\frac{6\sqrt{205}}{205}$.

 Applying \href{https://artofproblemsolving.com/wiki/index.php?title=Law\_of\_Sines}{Law of Sines} on $\bigtriangleup ABC$ we have
 \[\frac{AC}{\sin B}=\frac{AB}{\sin C}\]Solving gets $\sin B=\frac{56}{65}$.

 Applying \href{https://artofproblemsolving.com/wiki/index.php?title=Stewart\%27s\_Theorem}{Stewart's Theorem} with \href{https://artofproblemsolving.com/wiki/index.php?title=Cevian}{cevian} $AE$ we have
 \[(BC)(AE)^2+(BC)(BE)(CE)=(CE)(AB)^2+(BE)(AC)^2\]
 \[(BC)(AE)^2+(BC)(BE)(BC-BE)=(BC-BE)(AB)^2+(BE)(AC)^2\]Solving gets $AE=\sqrt{\frac{15BE^2-198BE+2535}{15}}$

 Finally, applying \href{https://artofproblemsolving.com/wiki/index.php?title=Law\_of\_Sines}{Law of Sines} on $\bigtriangleup BAE$ we have
 \[\frac{AE}{\sin B}=\frac{BE}{\sin BAE}\]
 \[\frac{\sqrt{\frac{15BE^2-198BE+2535}{15}}}{\frac{56}{65}}=\frac{BE}{\frac{6\sqrt{205}}{205}}\]
 \[7605BE^2-32342BE+2535=0\]Solving the easy quadratic equation gets $BE=\frac{1632}{463}\Longrightarrow q=\boxed{463}$

 





\textbf{Solution 10}

Making perpendicular lines from $D$ to $AC$, meeting at $N$; from $E$ to $AB$, meeting at $J$. According to LOC, we can get that $\cos\angle{C}=\frac{3}{5}$. So we get that $CN=\frac{18}{5};DN=\frac{24}{5};AN=AC-CN=\frac{52}{5}$. Now we can see that $\tan\angle{DAN}=\frac{DN}{AN}=\frac{6}{13}$. Now we see that in $\triangle{EJA}$, assuming $EJ=6x;AJ=13x$ since $\tan\angle{JAE}=\tan\angle{DAN}$. Now we need to find the tangent of $\angle{B}$. Making a perpendicular line from $C$ to $AB$ at $M$. We can see that $CM=\frac{168}{13};AM=\frac{70}{13};BM=AB-AM=13-\frac{70}{13}=\frac{99}{13}$. We get that $\tan\angle{B}=\frac{56}{33}$ so $BJ=\frac{33}{56}*6x=\frac{99x}{28}$. After getting AJ and BJ, we can get that $AB=AJ+BJ=\frac{463x}{28}=13$, which means that $x=\frac{364}{463}$. According to the similarity, $\frac{CM}{BC}=\frac{JE}{BE};BE=\frac{2535}{463}$ which $\boxed{463}$ is our answer~bluesoul

 





\textbf{Solution 11 (Ultimate Stewarts Bash)}

First, apply Stewart's theorem to triangle $ABC$ with cevian $AD$, from which we receive $AD=\frac{4\sqrt{41}}{\sqrt{5}}$. Then, set $AE=y$ and $BE=x.$ Hence, $DE=9-x$. Applying Stewart's on triangle $ACE$, with cevian $AD$, we receive that $y^2-x^2=169-\frac{66}{5}x$. By also applying the sine ratio formula on triangles $ACD$ and $AEB$, since these triangles share the same height, we get that $\frac{6}{x}=\frac{\frac{1}{2} \cdot 14 \cdot \frac{4\sqrt{41}}{\sqrt{5}}}{\frac{1}{2} \cdot 13 \cdot y}$. Here, we can simplify to receive that $\frac{32144}{7605}x^2=y^2$. We plug this into our earlier equation, and get that $\frac{24539}{7605}x^2=169-\frac{66}{5}x \implies 24539x^2+66 \cdot 39^2x-169 - 169 \cdot 39^2 \cdot 5=0.$ We then apply the quadratic formula (which may seem computationally intensive, but factoring kills it), and get $\frac{-66 \cdot 39^2 + \sqrt{66^2 \cdot 39^4+4(24539)(169 \cdot39^2 \cdot 5)}}{2 \cdot 24539}$. (Note we only take the plus symbol instead of $\pm$ since $x > 0$.) After factoring heavily, we get the answer to be $\frac{-33 \cdot 39^2+39 \cdot 13 \cdot 364}{24539} = \frac{2535}{463}$, and the answer is $\boxed{463}.$

 





\textbf{Solution 12 (Clean)}


                    \begin{figure}[H]
                        \centering
                        \includegraphics[width=0.5\textwidth]{Solution/AIME_2005_II_14_Solution_Figure_1.png}
                    \end{figure}
                    We begin by drawing a line through point $B$ parallel to side $AC$. Extend Lines $AE$ and $AD$ to meet the new parallel line at points $F$ and $G$, respectively. This will help us create a lot of similar triangles. 

 The first pair of similar triangles that is easy to spot is triangle $ACD$ and triangle $GBD$ by AA similarity ($\angle{ADC}=\angle{GDB}$ by vertical angles, and $\angle{DAC}=\angle{DGB}$ by alternate interior angles). With the ratios $\frac{AC}{BG}=\frac{CD}{BD}$ $\rightarrow$ $\frac{14}{BG}=\frac{6}{9}$ $\rightarrow$ $BG=21$. 

 The next pair of similar triangles is triangle $BAG$ and triangle $BFA$ by AA similarity ($\angle{CAD}=\angle{BAF}$ as given, and $\angle{CAD}=\angle{BGA}$ by alternate interior angles, so by transitive property, $\angle{BAF}=\angle{BGA}$. Additionaly, $\angle{ABF}=\angle{GBA}$ as they are the same angle.). With the ratios $\frac{AB}{BG}=\frac{BF}{AB}$ $\rightarrow$ $\frac{13}{21}=\frac{BF}{13}$ $\rightarrow$ $BF=\frac{169}{21}$. 

 The last pair of similar triangles we need is triangle $BEF$ and triangle $CAE$ by AA similarity (similar by vertical angles and alternate interior angles). With the ratios $\frac{AC}{BF}=\frac{CE}{BE}$ $\rightarrow$ $\frac{14}{\frac{169}{21}}=\frac{15-x}{x}$ $\rightarrow$ $x=\frac{169*15}{463}$. This cannot be simplified further, and the problem asks for the denominator, so the answer is $\boxed{463}$. 

 




	\item (\textit{AIME II 2005}) Let $w_1$ and $w_2$ denote the circles $x^2+y^2+10x-24y-87=0$ and $x^2 +y^2-10x-24y+153=0,$ respectively. Let $m$ be the smallest positive value of $a$ for which the line $y=ax$ contains the center of a circle that is externally tangent to $w_2$ and internally tangent to $w_1.$ Given that $m^2=\frac pq,$ where $p$ and $q$ are relatively prime integers, find $p+q.$


 
 \textbf{Answer: }169



\textbf{Solution 1}

Rewrite the given equations as $(x+5)^2 + (y-12)^2 = 256$ and $(x-5)^2 + (y-12)^2 = 16$.

 Let $w_3$ have center $(x,y)$ and radius $r$.  Now, if two circles with radii $r_1$ and $r_2$ are externally tangent, then the distance between their centers is $r_1 + r_2$, and if they are internally tangent, it is $|r_1 - r_2|$.  So we have

 
 \begin{align*} r + 4 &= \sqrt{(x-5)^2 + (y-12)^2} \\ 16 - r &= \sqrt{(x+5)^2 + (y-12)^2} \end{align*}


 Solving for $r$ in both equations and setting them equal, then simplifying, yields

 
 \begin{align*} 20 - \sqrt{(x+5)^2 + (y-12)^2} &= \sqrt{(x-5)^2 + (y-12)^2} \\ 20+x &= 2\sqrt{(x+5)^2 + (y-12)^2} \end{align*}


 Squaring again and canceling yields $1 = \frac{x^2}{100} + \frac{(y-12)^2}{75}.$

 So the locus of points that can be the center of the circle with the desired properties is an \href{https://artofproblemsolving.com/wiki/index.php?title=Ellipse}{ellipse}.

 
                    \begin{figure}[H]
                        \centering
                        \includegraphics[width=0.5\textwidth]{Solution/AIME_2005_II_15_Solution_Figure_1.png}
                    \end{figure}
                    Since the center lies on the line $y = ax$, we substitute for $y$ and expand: 
 \[1 = \frac{x^2}{100} + \frac{(ax-12)^2}{75} \Longrightarrow (3+4a^2)x^2 - 96ax + 276 = 0.\]

 We want the value of $a$ that makes the line $y=ax$ tangent to the ellipse, which will mean that for that choice of $a$ there is only one solution to the most recent equation. But a quadratic has one solution \href{https://artofproblemsolving.com/wiki/index.php?title=Iff}{iff} its discriminant is $0$, so $(-96a)^2 - 4(3+4a^2)(276) = 0$.

 Solving yields $a^2 = \frac{69}{100}$, so the answer is $\boxed{169}$.

 





\textbf{Solution 2}

As above, we rewrite the equations as $(x+5)^2 + (y-12)^2 = 256$ and $(x-5)^2 + (y-12)^2 = 16$. Let $F_1=(-5,12)$ and $F_2=(5,12)$. If a circle with center $C=(a,b)$ and radius $r$ is externally tangent to $w_2$ and internally tangent to $w_1$, then $CF_1=16-r$ and $CF_2=4+r$. Therefore, $CF_1+CF_2=20$. In particular, the locus of points $C$ that can be centers of circles must be an ellipse with foci $F_1$ and $F_2$ and major axis $20$.

 Clearly, the minimum value of the slope $a$ will occur when the line $y=ax$ is tangent to this ellipse. Suppose that this point of tangency is denoted by $T$, and the line $y=ax$ is denoted by $\ell$. Then we reflect the ellipse over $\ell$ to a new ellipse with foci $F_1'$ and $F_2'$ as shown below.

 
                    \begin{figure}[H]
                        \centering
                        \includegraphics[width=0.5\textwidth]{Solution/AIME_2005_II_15_Solution_Figure_1.png}
                    \end{figure}
                    By the reflection property of ellipses (i.e., the angle of incidence to a tangent line is equal to the angle of reflection for any path that travels between the foci), we know that $F_1$, $T$, and $F_2'$ are collinear, and similarly, $F_2$, $T$ and $F_1'$ are collinear. Therefore, $OF_1F_2F_2'F_1'$ is a pentagon with $OF_1=OF_2=OF_1'=OF_2'=13$, $F_1F_2=F_1'F_2'=10$, and $F_1F_2'=F_1'F_2=20$. Note that $\ell$ bisects $\angle F_1'OF_1$. We can bisect this angle by bisecting $\angle F_1'OF_2$ and $F_2OF_1$ separately.

 We proceed using complex numbers. Triangle $F_2OF_1'$ is isosceles with side lengths $13,13,20$. The height of this from the base of $20$ is $\sqrt{69}$. Therefore, the complex number $\sqrt{69}+10i$ represents the bisection of $\angle F_1'OF_2$.

 Similarly, using the 5-12-13 triangles, we easily see that $12+5i$ represents the bisection of the angle $F_2OF_1$. Therefore, we can add these two angles together by multiplying the complex numbers, finding
 \[\text{arg}\left((\sqrt{69}+10i)(12+5i)\right)=\frac{1}{2}\angle F_1'OF_1.\]Now the point $F_1$ is given by the complex number $-5+12i$. Therefore, to find a point on line $\ell$, we simply subtract $\frac{1}{2}\angle F_1'OF_1$, which is the same as multiplying $-5+12i$ by the conjugate of $(\sqrt{69}+10i)(12+5i)$. We find
 \[(-5+12i)(\sqrt{69}-10i)(12-5i)=169(10+i\sqrt{69}).\]In particular, note that the tangent of the argument of this complex number is $\sqrt{69}/10$, which must be the slope of the tangent line. Hence $a^2=69/100$, and the answer is $\boxed{169}$.

 





\textbf{Solution 3}

We use the same reflection as in Solution 2. As $OF_1'=OF_2=13$, we know that $\triangle OF_1'F_2$ is isosceles. Hence $\angle F_2F_1'O=\angle F_1'F_2O$. But by symmetry, we also know that $\angle OF_1T=\angle F_2F_1'O$. Hence $\angle OF_1T=\angle F_1'F_2O$. In particular, as $\angle OF_1T=\angle OF_2T$, this implies that $O, F_1, F_2$, and $T$ are concyclic.

 Let $X$ be the intersection of $F_2F_1'$ with the $x$-axis. As $F_1F_2$ is parallel to the $x$-axis, we know that 
 \[\angle TXO=180-\angle F_1F_2T.\tag{1}\] But 
 \[180-\angle F_1F_2T=\angle F_2F_1T+\angle F_1TF_2.\tag{2}\] By the fact that $OF_1F_2T$ is cyclic, 
 \[\angle F_2F_1T=\angle F_2OT\qquad\text{and}\qquad \angle F_1TF_2=\angle F_1OF_2.\tag{3}\] Therefore, combining (1), (2), and (3), we find that
 \[\angle TXO=\angle F_2OT+\angle F_1OF_2=\angle F_1OT.\tag{4}\]

 By symmetry, we also know that
 \[\angle F_1TO=\angle OTF_1'.\tag{5}\]Therefore, (4) and (5) show by AA similarity that $\triangle F_1OT\sim \triangle OXT$. Therefore, $\angle XOT=\angle OF_1T$.

 Now as $OF_1=OF_2'=13$, we know that $\triangle OF_1F_2'$ is isosceles, and as $F_1F_2'=20$, we can drop an altitude to $F_1F_2'$ to easily find that $\tan \angle OF_1T=\sqrt{69}/10$. Therefore, $\tan\angle XOT$, which is the desired slope, must also be $\sqrt{69}/10$. As before, we conclude that the answer is $\boxed{169}$.

 





\textbf{Solution 4}


                    \begin{figure}[H]
                        \centering
                        \includegraphics[width=0.5\textwidth]{Solution/AIME_2005_II_15_Solution_Figure_1.png}
                    \end{figure}
                    First, rewrite the equations for the circles as $(x+5)^2+(y-12)^2=16^2$ and $(x-5)^2+(y-12)^2=4^2$.  Then, choose a point $(a,b)$ that is a distance of $x$ from both circles.  Use the distance formula between $(a,b)$ and each of $A$ and $C$ (in the diagram above).  The distances, as can be seen in the diagram above are $16-x$ and $4+x$, respectively.
 \[(a-5)^2+(b-12)^2=(4+x)^2\]
 \[(a+5)^2+(b-12)^2=(16-x)^2\]Subtracting the first equation from the second gives 
 \[20a=240-40x\rightarrow a=12-2x\rightarrow x=6-\frac a2\]Substituting this into the first equation gives
 \[a^2-10a+25+b^2-24b+144=100-10a+\frac{a^2}4\]
 \[b^2-24b+69+\frac{3a^2}4=0\]Now, instead of converting this to the equation of an eclipse, solve for $b$ and then divide by $a$.
 \[b=\frac{24\pm\sqrt{300-3a^2}}{2}\]We take the smaller root to minimize $\frac b a$.
 \[\frac b a=\frac{24-\sqrt{300-3a^2}}{2a}=\frac{24-\sqrt3\cdot\sqrt{100-a^2}}{2a}=\frac{12}a-\frac{\sqrt3}{2a}\sqrt{100-a^2}\]Now, let $10\cos\theta=a$.  This way, $\sqrt{100-a^2}=10\sin\theta$.Substitute this in.  $\frac{b}{a}=\frac{12}{10\cos\theta}-\frac{\sqrt3\cdot\sin\theta}{2\cos\theta}=\frac65\sec\theta-\frac{\sqrt3}{2}\tan\theta$Then, take the derivative of this and set it to 0 to find the minimum value.$\frac{6}{5}\sec\theta\tan\theta-\frac{\sqrt3}{2}\sec^2\theta=0\rightarrow\frac{6}{5}\sin\theta-\frac{\sqrt3}{2}=0\rightarrow\sin\theta=\frac{5\sqrt3}{12}$Then, use this value of $\sin\theta$ to find the minimum of $\frac65\sec\theta-\frac{\sqrt3}{2}\tan\theta$ to get $\frac{\sqrt{69}}{10}\rightarrow\left(\frac{\sqrt{69}}{10}\right)^2=\frac{69}{100}\rightarrow69+100=\boxed{169}$

 





\textbf{Solution 5 (probably fastest)}

Like before, notice that the distances from the centers of the given circles to the desired center are $4+r$ and $16-r$, which add up to $20$. This means that the possible centers of the third circle lie on an ellipse with foci $(-5, 12)$ and $(5, 12)$. Using the fact that the sum of the distances from the foci is $20$, we find that the semi-major axis has length $10$ and the semi-minor axis has length $5\sqrt{3}$. Therefore, the equation of the ellipse is 
 \[\dfrac{x^2}{100}+\dfrac{(y-12)^2}{75} = 1,\] where the numbers $100$ and $75$ come from $10^2$ and $(5\sqrt{3})^2$ respectively. 

 
We proceed to find $m$ using the same method as Solution 1.

 





\textbf{Solution 6}

First, obtain the equation of the ellipse as laid out in previous solutions. We now scale the coordinate plane in the $x$ direction by a factor of $\frac{\sqrt{3}}{2}$ centered at $x=0.$ This takes the ellipse to a circle centered at $(0,12)$ with radius $5\sqrt{3}$ and takes the line $y=ax$ to $y=\left( \frac{\sqrt{3}}{2} \right)^{-1} ax.$ The tangent point of our line to the circle with positive slope forms a right triangle with the origin and the center of the circle. Thus, the distance from this tangent point to the origin is $\sqrt{69}.$ By similar triangles, the slope of this line is then $\frac{\sqrt{69}}{5\sqrt{3}}.$ We multiply this by $\frac{\sqrt{3}}{2}$ to get $a=\frac{\sqrt{69}}{10},$ so our final answer is $\boxed{169.}$

 




\end{enumerate}